\section{Worked examples}

\subsection{\texorpdfstring{A successful \(\forall\exists\forall\) challenge}{A successful forall-exists-forall challenge}}

Consider
\[
 A\equiv
 \forall q<4\;\exists z<4\;\forall e<3\;
     ((z+q)\bmod 4=0).
\]
The variable \(e\) is deliberately irrelevant to the atom; it forces a
nontrivial universal certificate table and makes the code structure visible.
For every \(q\), choose \(z=(-q)\bmod4\).

At an atom, both code bounds are one.  For the inner universal over three
values of \(e\), the positive bound is \(1^3=1\), while the negative bound is
\(1+1+1=3\).  For the existential over four values of \(z\),
\[
 P^+=1+1+1+1=4,
 \qquad
 P^-=3^4=81.
\]
For the outer universal over four values of \(q\),
\[
 P_A^+=4^4=256,
 \qquad
 P_A^-=81+81+81+81=324.
\]
The reference implementation computes value one and canonical positive code
108.  In base four, the four outer table digits encode the least successful
\(z\) for \(q=0,1,2,3\), namely \(0,3,2,1\) under the chosen little-endian
convention.  Each digit points to the unique inner positive table code zero.

The positive checker decodes all four digits.  At each \(q\) it locates the
selected \(z\), checks the three atomic leaves, and introduces the existential.
Universal introduction then assembles the four verified responses.  A malformed
code can be rejected at a bound test, at an owner interval, or at an atomic
recomputation; the proof trusts none of the producer's intentions.

The generated fixture is
\code{examples/bounded\_challenge\_success.json}.  It records the formula,
code bounds, canonical code, and a fully decoded human-readable certificate
tree.

\subsection{A least counterchallenge}

Now consider
\[
 B\equiv
 \forall q<3\;\exists z<2\;\forall e<2\;(q<z).
\]
For \(q=0\), the witness \(z=1\) works.  For \(q=1\), neither \(z=0\) nor
\(z=1\) satisfies \(1<z\).  Thus the least failing outer challenge is
\(q=1\).

The inner universal has positive bound one and negative bound two.  The
existential has positive bound two and negative bound \(2^2=4\).  The outer
universal therefore has
\[
 P_B^+=2^3=8,
 \qquad
 P_B^-=4+4+4=12.
\]
The canonical negative code is 4: the first interval \([0,4)\) belongs to
\(q=0\), and the next interval begins at 4 and belongs to \(q=1\).  Its local
code is zero, a mixed-radix table containing negative certificates for both
candidate witnesses.  Each candidate's negative universal certificate chooses
the least failing \(e\), here \(e=0\).

Negative checker soundness proceeds exactly as the intended counterargument.
Assume the outer universal.  Instantiate it at \(q=1\).  Open a proposed
\(z<2\).  Decode the corresponding table digit, obtain a negative certificate
for the inner universal, and contradict that universal.  Therefore \(\neg B\).
The fixture
\code{examples/bounded\_challenge\_failure.json} records this counterstrategy.

\subsection{Empty bounds}

Let
\[
 E_0=\exists x<0\;\top,
 \qquad
 U_0=\forall x<0\;\bot.
\]
The value of \(E_0\) is zero; its positive code bound is the empty sum zero,
and its negative code bound is the empty product one.  Code zero is the empty
negative table.  Conversely \(U_0\) has positive bound one, negative bound zero,
and code zero is its positive empty table.

These cases are important regression tests.  A sentinel-based implementation
that confuses the bound with a valid witness often fails at zero.  The
disjoint-sum/product definitions handle them without exceptions.  The exact
outputs are in \code{examples/empty\_bounds.json}.

\subsection{A cyclic demand closure}

Take capacities \((3,3,4)\), seed \((0,0,0)\), and monotone requirements
\begin{align*}
 F_0(d)&=\max(1,d_1),\\
 F_1(d)&=\min(3,d_0+1),\\
 F_2(d)&=\max(2,d_0,d_1).
\end{align*}
Iteration of \(d\mapsto d\vee F(d)\) gives
\[
 (0,0,0),
 (1,1,2),
 (1,2,2),
 (2,2,2),
 (2,3,2),
 (3,3,3).
\]
The final vector is stable and is the least prefixed point above the seed.  The
rank budget is \(3+3+4=10\), while only five strict steps are needed.  This
example demonstrates why a cyclic graph need not be broken by guessing a large
global precision.  The exact trace appears in
\code{examples/cyclic\_demand\_closure.json}.

\subsection{Finite maximum on a certified net}

Suppose a domain certificate lists \(y_0,\ldots,y_{N-1}\) and proves that every
admissible \(x\) lies within \(\delta\) of some \(y_i\).  Suppose each value
\(f(y_i)\) is approximated by a rational \(r_i\) within \(\eta\), and a modulus
proves \(d(x,y)\le\delta\to |f(x)-f(y)|\le\varepsilon\).

A bounded formula compares the rational codes and chooses the least index
\(i^*\) maximizing \(r_i\).  Its positive strategy certifies every comparison
used by the finite argmax.  For arbitrary admissible \(x\), choose a covered
\(y_j\).  Then
\begin{align*}
 f(x)&\le f(y_j)+\varepsilon\\
     &\le r_j+\eta+\varepsilon\\
     &\le r_{i^*}+\eta+\varepsilon\\
     &\le f(y_{i^*})+2\eta+\varepsilon.
\end{align*}
Thus the selected sample is within \(2\eta+\varepsilon\) of the supremal value.
The cover and modulus are explicit residual producers.  The finite comparison,
argmax, and failure explanation are generic STRIVE components.

\subsection{Approximate reciprocal}

Let rational observations satisfy
\[
 |x-\tilde x|\le\delta,
 \qquad |x|\ge\lambda>0,
 \qquad \delta\le\lambda/2.
\]
Then \(|\tilde x|\ge\lambda/2\), and
\[
 \left|\frac1x-\frac1{\tilde x}\right|
 =\frac{|x-\tilde x|}{|x|\,|\tilde x|}
 \le \frac{2\delta}{\lambda^2}.
\]
All operations can be expressed by cross-multiplied integer inequalities.  A
reverse demand transformer chooses \(\delta\) so that
\(2\delta/\lambda^2\le2^{-n}\).  The lower bound \(\lambda\) is the crucial
producer; \(x\ne0\) alone is not substituted for it.

\subsection{Bounded Diophantine completeness}

Let \(p(\vec x)=0\) be a fixed polynomial and suppose a residual theorem gives
\(|x_i|\le H\) for every solution.  Encode signed integers by naturals and
enumerate the finite box.  A positive certificate for ``\(\vec x\) is a
solution'' is just the arithmetic evaluation trace.  A negative table proves
that every nonlisted tuple fails.  The output list is complete relative to the
height theorem.

This example isolates the genuine bottleneck.  The strategy compiler solves
bounded enumeration and completeness checking.  It does not derive \(H\) from
a statement that the solution set is finite.
